\documentclass[a4paper]{article}

\usepackage[english]{babel}
%\usepackage[utf8x]{inputenc}
\usepackage{amsmath}
\usepackage{graphicx}
\usepackage[colorinlistoftodos]{todonotes}
\usepackage[all]{xy}
\usepackage{dsfont}
\usepackage{amssymb}
\usepackage{bbold}
\usepackage{amsthm}

\usepackage{tikz}
\usepackage{verbatim}
\usetikzlibrary{matrix}
\usetikzlibrary{cd}

\usepackage[autostyle=true]{csquotes}
\usepackage{pdfpages}
\usepackage{hyperref}
\usepackage{fancyhdr}


\newtheorem{theorem}{Theorem}[section]
\newtheorem{proposition}{Proposition}[theorem]
\newtheorem{lemma}{Lemma}[theorem]
\newtheorem{corollary}{Corollary}[theorem]
\newtheorem*{definition}{Definition}


\newcommand{\norm}[1]{\left\lVert#1\right\rVert}

\newcommand{\Rn}{\mathbb{R}^n}


\newcommand{\too}[1]{\overset{#1}{\longrightarrow}}
\newcommand{\hookar}[1]{\overset{#1}{\hookrightarrow}}
\newcommand{\colim}[1]{ \underset{#1}{\texttt{colim}} \text{ } }
\newcommand{\fset}[1]{ \{ #1_0 , \dots , #1_n \} }
\newcommand{\ddt}{\frac{d}{dt}}


\newcommand{\N}{\mathbb{N}}
\newcommand{\Z}{\mathbb{Z}}





\title{Laboratorio 1}
\date{}
\fancyhead[L]{Universidad Finis Terrae \\ Facultad de Ingeniería \\ Algebra Lineal \\ 2do semestre de 2023}
\fancyhead[R]{\includegraphics[scale=0.1]{UFT.png}}

\begin{document}

\maketitle
\thispagestyle{fancy}

\section*{Los Enteros módulo $n$}
\vspace{0.2 cm}

Los enteros mód $n$ (see lee m\'odulo $n$) se han vuelto indispensables en la teoría y las aplicaciones del álgebra. Por ejemplo, se pueden aplicar en criptografía, teoría de códigos, y la detección de errores en códigos de identificación.

Dos enteros $a$ y $b$ son equivalentes mód $n$, es decir, representan el mismo elemento, si 
$n$ divide a $a-b$.

Los enteros mód $n$ también particionan 
$\mathbb{Z}$ en $n$ distintos subconjuntos; denotaremos el conjunto de subconjuntos(llamados clases de equivalencia) por $\mathbb{Z}_n$.

En \textit{Mathlab} podemos escribir

\[\texttt{b=mod(a,n)}\]
entrega el resto después de dividir $a$ por $n$, donde $a$ es el dividendo y $n$ es el divisor. La función \texttt{mod} sigue la convención de que \texttt{mod(a,0)} devuelve $a$.

% Considere los enteros módulo 12 y la correspondiente partición de los enteros, en las siguientes clases de equivalencia:
%\begin{align*}
%    [0]&=\{\dots, -12,0,12,24,\dots\}\\
%    [1]&=\{\dots, -11,0,13,25,\dots\}\\
%    \vdots & \\
%    [11]&=\{\dots, -1,11,23,35,\dots\} 
%\end{align*}
%Por ejemplo, escribimos que la clase de equivalencia de cero, es igual a la de 12
%\[[0]=[12]\]
%Cuando no haya posibilidad de confusión, usaremos $0,1,\dots,11$ para indicar las clases de equivalencia $$[0],[1],\dots,[11]$$ respectivamente. 
Podemos hacer aritmética en $\mathbb{Z}_n$. Para dos enteros $a$ y $b$, definimos adición módulo $n$ como $(a+b)\mod n$; es decir, el resto de la división de $a+b$ entre $n$. Similarmente, la multiplicación módulo $n$ se define como $(ab)\mod n$, el resto de la división de $ab$ entre $n$. Por ejemplo, en $\mathbb{Z}_{12}$, la suma $[5]+[9]=[14]=[2]$.

M\'as a\'un, nos interesan los vectores con coordenadas en $\mathbb{Z}_n$. Por ejemplo, $\mathbb{Z}^m_n=\{(a_1,a_2,\dots,a_m)\;:\; a_i \in \mathbb{Z}_n\}$. Defina una operación en $\mathbb{Z}^m_n$ por
\[(a_1,a_2,\dots,a_m)+(b_1,b_2,\dots,b_m)=(a_1+b_1,a_2+b_2,\dots,a_m+b_m),\]
es decir, la suma usual de vectores. Luego se calcula a qu\'e clase de equivalencia corresponde cada coordenada.


\section*{N\'umeros binarios}
\vspace{0.2 cm}

Usualmente, usamos el sistema numérico \textbf{decimal} para representar n\'umeros, usando dígitos entre 0 y 9. Escribimos diferentes cantidades utilizando una combinación de estos símbolos y cada d\'igito adicional aumenta el orden de magnitud en 10. Por ejemplo, $65=6\cdot 10+5$. Cuando escribamos en sistema decimal, usaremos la notaci\'on $65_{10}$.

En el sistema \textbf{binario}, los números se representan usando s\'olo dos símbolos diferentes que son el 0 y el 1. En este caso, cada posici\'on (de derecha a izquierda) nos indica la potencia de 2 que se sumar\'a. Siguiendo con el ejemplo anterior, \[65_{10}=1000001=2^6+0+0+0+0+0+2^0=64+1.\] Cuando escribamos en sistema binario, usaremos la notaci\'on $1000001_{2}$.
\begin{itemize}

\end{itemize}

\section*{Teor\'ia de c\'odigos}
\vspace{0.2 cm}
En teor\'ia de c\'odigos los mensajes se transforman en vectores en $\Z_2^n$, donde cada coordenada se conoce como \emph{bit}.  El conjunto de todos los vectores posible se llama código y la transformaci\'on se hace a través de un mecanismo llamado codificador. El mensaje es transmitido y luego decodificado. 

Consideraremos la aparición de errores durante la transmisión. Un error ocurre si hay un cambio en uno o más bits de la palabra del código. Un protocolo decodificador es un método que ya sea convierte 
vectores en $\Z_2^n$ arbitraria recibida en un mensaje decodificado coherente o da un mensaje de error para ese vector. 

Si el mensaje recibido es una palabra del código (vector permitido), entonces el mensaje decodificado debe ser el mensaje que fue codificado en la palabra del código. 

Para los vectores recibidos que no están en el código, el protocolo dará una indicación de error, o, si somos más astutos, tratará de corregir el error y reconstruir el mensaje original. Nuestro objetivo es transmitir mensajes libres de errores de la forma más barata y rápida posible.

\vspace{0.2 cm}
\textbf{Problema a resolver:}
\vspace{0.2 cm}





\begin{enumerate}
\item Estudie el conjunto $\mathbb{Z}_2$ y haga una tabla de las operaciones suma y multiplicaci\'on. Verifique usando \emph{Mathlab}.
\item Investigue sobre el código ASCII (American Standard Code for Information Interchange) como vectores en $\Z_2^n$, explicitando $n$, el n\'umero total de vectores posibles, n\'umero de caracteres ASCII, etc).
\item En el código ASCII, ¿qué se puede o debe hacer con el/los bits restantes? Investigue sobre \emph{la paridad} en el c\'odigo y como se relaciona con el producto punto de un vector $x=(x_1,\dots, x_n)$ que todos los d\'igitos son 1, es decir,
\[\textbf{1}^t\cdot x=1\cdot x_1+\cdots+ 1\cdot x_n.\]
¿Cómo se relaciona la paridad con la detecci\'on de errores? ¿En un mismo mensaje se puede detectar m\'as de un error? Justifique.
\item Sea $H$ una matriz en $\mathcal{M}_{m \times n}(\Z^2)$, es decir, cuyas entradas
son 0 o 1. Se denota como espacio nulo de $H$ al conjunto
\[\operatorname{Nul}(H)=\{y\in \Z^n_2\colon Hy=0\},\]
dicho conjunto se conoce como c\'odigo lineal.
Verifique para la matriz 
\[H=\begin{pmatrix}
   0& 1& 0& 1& 0 \\
    1& 1& 1& 1& 0 \\
     0& 0& 1& 1& 1 
\end{pmatrix}\]
si los vectores $(1 1 1 1 0 0), (01011),(1111110), (010101)$ pertenecen o no al c\'odigo. En \emph{Mathlab}, use el comando \texttt{null}$(H)$ para encontrar los vectores en el c\'odigo.
\end{enumerate}


\textbf{Aspectos formales del Laboratorio}
{\bf{Grupos de Trabajo}:} Los grupos de trabajo estarán conformados por 3 o 4 estudiantes y
deben ser aprobados por el profesor. Se penalizarán las modificaciones posteriores de
estos grupos.
\vspace*{.5cm}

Todas las entregas se deben realizar en el portal e-Finis en los buzones habilitados para
tal efecto, no se aceptan trabajos enviados por correo. 

\vspace*{.5cm}

\noindent{\bf{Formato de Entrega}}
Se deben entregar los documentos en PDF (en caso contrario no se corregirá el informe)
en formato \LaTeX o con el formato entregado por su profesor.



{\bf{NOTA}}
La copia y/o plagio serán sancionados de acuerdo a lo estipulado en el Reglamento de la
Universidad.
\end{document}